\documentclass[11pt]{article}
\topmargin=-1.0cm
\usepackage{graphicx}
\usepackage{epstopdf}
\usepackage{amssymb}
\usepackage{amsmath}
\usepackage{amsthm}
\usepackage{color}
\usepackage[margin=1in]{geometry}
\usepackage[hidelinks]{hyperref}
\usepackage{booktabs}

% ******************************* Author Macros *******************************

\newcommand{\pl}[2]{\frac{\partial#1}{\partial#2}}
\newcommand{\bu}{{\bf u} }
\newcommand{\p}{\partial}
\newcommand{\og}{\omega}
\newcommand{\Og}{\Omega}
\newcommand{\fl}[2]{\frac{#1}{#2}}
\newcommand{\dt}{\delta}
\newcommand{\nn}{\nonumber}
\newcommand{\ap}{\alpha}
\newcommand{\bt}{\beta}
\newcommand{\tht}{\theta}
\newcommand{\kp}{\kappa}
\newcommand{\Dt}{\Delta}
\renewcommand{\theequation}{\arabic{section}.\arabic{equation}}
\newcommand{\be}{\begin{equation}}
\newcommand{\ee}{\end{equation}}
\newcommand{\ba}{\begin{array}}
\newcommand{\ea}{\end{array}}
\newcommand{\bea}{\begin{eqnarray}}
\newcommand{\eea}{\end{eqnarray}}
\newcommand{\beas}{\begin{eqnarray*}}
\newcommand{\eeas}{\end{eqnarray*}}
\newtheorem{theorem}{Theorem}[section]
\newtheorem{lemma}{Lemma}[section]
\newtheorem{proposition}{Proposition}[section]
\newtheorem{remark}{Remark}[section]
\newtheorem{assumption}{Assumption}[section]
\newcommand{\bx}{{\bf x} }
\newcommand{\gm}{\gamma}
\newcommand{\vep}{\varepsilon}
\newcommand{\T}{\mathbb T}
\newcommand{\dT}{d_{\mathbb T}}

% *****************************************************************************

\title{A WKB-based numerical method for capturing trait concentration in a dispersal evolution model}

\author{
Weijie Huang\thanks{School of Mathematics and Statistics, Beijing Jiaotong University, Beijing 100044, People's Republic of China {\tt (wjhuang@bjtu.edu.cn)}} \ and
Xinran Ruan\thanks{School of Mathematical Sciences, Capital Normal University, Beijing, China, 100048, People's Republic of China {\tt (xinran.ruan@cnu.edu.cn)}.}
}

\begin{document}
\date{}
\maketitle

\begin{abstract}
The evolution of a dispersal trait is a classical question in evolutionary ecology. A typical phenomenon is that, in the rare mutation regime, the trait distribution concentrates toward a Dirac mass, which makes the accurate numerical capture of the fittest trait difficult. In this paper, we consider a WKB reformulation of a dispersal evolution model and propose a semi-discrete numerical method for resolving the concentration in the trait direction without using prohibitively fine trait grids. We present the detailed semi-discrete scheme, establish positivity of the amplitude variable, and derive a fixed-grid stationary AP structure under an explicit one-sided stability condition on the weighted WKB remainder. We also identify a density-compatible WKB variant for which this remainder condition follows from discrete summation by parts. These results clarify both the useful structure and the current limitations of the method. Numerical experiments are organized to compare the proposed WKB formulation with direct discretizations, to monitor the weighted remainder, and to illustrate its advantage in locating the fittest trait in the small-mutation regime.
\end{abstract}

{\bf Key words. } dispersal evolution, trait concentration, WKB ansatz, asymptotic-preserving method, Hamilton--Jacobi equation, principal eigenvalue.


%===================================================================================================
%	Introduction
%===================================================================================================
\section{Introduction}

Trait concentration in the small-mutation regime is a major source of difficulty in the numerical approximation of space--trait structured population models. When the mutation parameter is small, the solution typically develops a sharp peak in the trait variable, while the dominant trait is selected through a delicate interaction between spatial heterogeneity, dispersal, and competition. In this regime, a direct discretization of the original density often requires a prohibitively fine mesh in the trait direction in order to resolve the narrow peak and identify the fittest trait accurately. As a consequence, the main numerical issue is not merely to approximate the full density, but to capture the location and structure of the emerging concentration efficiently and robustly.

In this work, we consider a representative dispersal evolution model of the form
\begin{equation}\label{eq:n_intro}
    \varepsilon \partial_t n_\varepsilon
    -D(\theta)\Delta_\mathbf{x} n_\varepsilon
    -\varepsilon^2\Delta_\theta n_\varepsilon
    =
    n_\varepsilon\bigl(K(\mathbf{x})-\rho_\varepsilon(\mathbf{x})\bigr),
    \qquad \mathbf{x}\in\Omega,
\end{equation}
where $t$ denotes time, $\mathbf{x}$ the spatial variable, and $\theta$ a phenotypic variable. The unknown $n_\varepsilon(t,\mathbf{x},\theta)$ is the structured population density, and
\[
\rho_\varepsilon(t,\mathbf{x})
=
\int_{\mathbb T} n_\varepsilon(t,\mathbf{x},\theta)\,d\theta
\]
is the total population density. The function $K(\mathbf{x})$ represents the local carrying capacity, while the trait-dependent diffusivity $D(\theta)$ models the dispersal rate associated with phenotype $\theta$. 
Throughout the paper, we assume that the dispersal rate is positive and periodic in the trait variable and that the carrying capacity is nonnegative and bounded:
\begin{equation}\label{eq:basic_assumptions}
    0<D_{\min}\le D(\theta)\le D_{\max}<\infty,
    \qquad
    0\le K(\mathbf{x})\le K_{\max}.
\end{equation}
We impose homogeneous Neumann boundary conditions on $\partial\Omega$ and periodic boundary conditions in the $\theta$-direction. Although our numerical development is carried out for \eqref{eq:n_intro}, the main difficulty addressed here is more general and arises in a broader class of small-mutation trait-selection problems whose solutions become highly concentrated in the phenotypic direction.

Model \eqref{eq:n_intro} and related equations have been extensively studied from the analytical viewpoint, including well-posedness, long-time behavior, and the connection with adaptive dynamics \cite{MirrahimiPerthame,PerthameSouganidis,AlmeidaPerthameRuan}. A characteristic feature is that the competition between spatial heterogeneity and trait-dependent dispersal may select an optimal phenotype in the long-time regime. In many situations, the steady state exhibits concentration in the trait variable. The small parameter $\varepsilon>0$ plays the role of a rare mutation scale, and in the limit $\varepsilon\to0$ the steady state typically converges, in a weak sense, to a Dirac mass supported at the fittest trait. This singular limit is well understood at the PDE level through WKB or Hopf--Cole transforms and constrained Hamilton--Jacobi equations coupled with principal eigenvalue problems \cite{PerthameSouganidis,BarlesPerthame,MirrahimiPerthame}.

From the numerical point of view, however, this asymptotic regime remains challenging. As $\varepsilon$ decreases, the trait profile of $n_\varepsilon$ becomes increasingly narrow, so standard discretizations of the original variable require a rapidly growing number of grid points in the phenotypic direction. Moreover, approximating the full density with reasonable accuracy does not automatically guarantee an accurate identification of the dominant trait. This is particularly restrictive in the small-$\varepsilon$ regime, where the quantity of real interest is often the location of the concentration and its effective profile rather than a pointwise resolution of the entire density on an extremely fine mesh.

Motivated by the asymptotic structure of the rare-mutation limit, we adopt the WKB representation
\[
n_\varepsilon(\mathbf{x},\theta,t)
=
W_\varepsilon(\mathbf{x},\theta,t)\,
\exp\!\left(\frac{u_\varepsilon(\theta,t)}{\varepsilon}\right),
\]
where the phase variable $u_\varepsilon$ describes the rapidly varying exponential profile in the trait direction, while $W_\varepsilon$ is a comparatively smoother amplitude. In the long-time and rare-mutation regime, the effective behavior of the phase variable is related to a constrained Hamilton--Jacobi equation involving a principal eigenvalue problem in the spatial variable \cite{PerthameSouganidis,AlmeidaPerthameRuan}. This observation suggests that the WKB reformulation is not only asymptotically natural but also computationally advantageous. In particular, it provides a mechanism for extracting the dominant trait from the phase variable without directly resolving the narrow peak of the original density on a prohibitively fine trait grid.

The present paper is organized around this computational viewpoint. Our main goal is to develop and study a WKB-based numerical strategy that is able to capture trait concentration more efficiently than direct discretizations of the original model. To this end, we first derive a semi-discrete formulation for the phase and amplitude variables. We then specify a concrete numerical method based on a monotone discretization of the Hamilton--Jacobi part and a positivity-preserving discretization of the amplitude equation. At the analytical level, we establish positivity of the amplitude and derive a stationary, fixed-grid AP structure under transparent reconstruction and one-sided remainder stability assumptions. Since the WKB split does not automatically provide an exact discrete chain rule, we also record a density-compatible variant that restores a conservative density-level balance and makes the weighted remainder estimate verifiable. On the numerical side, our experiments are designed to compare the original discretization and the WKB formulation directly, to track the remainder diagnostic, and to test the accurate capture of the fittest trait when $\varepsilon$ is small.

The paper is organized as follows. In Section~2 we recall the WKB ansatz and the formal steady-state rare mutation limit that motivate the numerical formulation. In Section~3 we present the semi-discrete scheme, prove positivity, and derive the conditional stationary AP structure, including the density-compatible variant that controls the weighted remainder. Section~4 is devoted to numerical experiments, where we compare the WKB formulation with direct discretizations, examine its performance in the concentration regime, and monitor the remainder diagnostic. The appendices collect typical choices of monotone numerical Hamiltonians and monotone numerical fluxes, a fully discrete realization, and sufficient conditions for the one-sided remainder bound.


%===================================================================================================
%   Preliminary results
%===================================================================================================
\section{Preliminary}

\subsection{WKB ansatz and reformulated system}

To capture the concentration phenomenon in the trait variable, we introduce the classical WKB ansatz
\begin{equation}\label{eq:WKB}
   n_\varepsilon(x,\theta,t)
   = W_\varepsilon(x,\theta,t)\,
   \exp\!\left(\frac{u_\varepsilon(\theta,t)}{\varepsilon}\right).
\end{equation}
This decomposition separates the exponentially varying part in the trait direction from a comparatively smoother amplitude. It is therefore natural for both the continuous asymptotic analysis and the numerical design in the small-mutation regime.

Substituting \eqref{eq:WKB} into \eqref{eq:n_intro} and assigning the leading trait-direction exponential contributions to the phase function, we obtain the reformulated system
\begin{equation}\label{eq:WKB_reformulation0}
\left\{
\begin{aligned}
& \partial_t u_\varepsilon
- |\partial_\theta u_\varepsilon|^2
- \varepsilon \, \partial_{\theta\theta} u_\varepsilon
= - \mathcal H(\theta,t), \\
& \varepsilon \partial_t W_\varepsilon
- D(\theta)\Delta_x W_\varepsilon
- \varepsilon^2 \partial_{\theta\theta} W_\varepsilon
- 2\varepsilon \, \partial_\theta W_\varepsilon \, \partial_\theta u_\varepsilon
= W_\varepsilon\bigl(K(x)-\rho_\varepsilon(x,t)+\mathcal H(\theta,t)\bigr),
\end{aligned}
\right.
\end{equation}
where \(\mathcal H(\theta,t)\) denotes an effective Hamiltonian. 
Using the identity
\[
\partial_\theta(W_\varepsilon\partial_\theta u_\varepsilon)
=
\partial_\theta W_\varepsilon\,\partial_\theta u_\varepsilon
+
W_\varepsilon\partial_{\theta\theta}u_\varepsilon,
\]
we rewrite the transport coupling in conservative form. If
\begin{equation}\label{def:F}
F_\varepsilon=-W_\varepsilon\partial_\theta u_\varepsilon,
\end{equation}
then
\[
-2\varepsilon\partial_\theta W_\varepsilon\,\partial_\theta u_\varepsilon
=
2\varepsilon\partial_\theta F_\varepsilon
+
2\varepsilon W_\varepsilon\partial_{\theta\theta}u_\varepsilon .
\]
Consequently, the system can be written as
\begin{equation}\label{eq:WKB_reformulation}
\left\{
\begin{aligned}
& \partial_t u_\varepsilon
- |\partial_\theta u_\varepsilon|^2
- \varepsilon \, \partial_{\theta\theta} u_\varepsilon
= - \mathcal H(\theta,t), \\
& \varepsilon \partial_t W_\varepsilon
- D(\theta)\Delta_x W_\varepsilon
- \varepsilon^2 \partial_{\theta\theta} W_\varepsilon
+2\varepsilon\partial_\theta F_\varepsilon
=
W_\varepsilon
\bigl(
K(x)-\rho_\varepsilon(x,t)
+\mathcal H(\theta,t)
-2\varepsilon\partial_{\theta\theta}u_\varepsilon
\bigr).
\end{aligned}
\right.
\end{equation}
In the continuous theory, \(\mathcal H\) is related to a principal eigenvalue problem in the spatial variable. This structure is the basic reason for introducing the variables \((u_\varepsilon,W_\varepsilon)\) in the first place.

\subsection{Continuous asymptotic structure of the steady state}\label{subsec:continuous_asymptotic_structure}

We next recall the steady-state asymptotic structure that motivates the numerical method. Let \((\bar n_\varepsilon,\bar \rho_\varepsilon)\) be a positive steady state of \eqref{eq:n_intro}. Then
\begin{equation}\label{eq:ss_n}
-D(\theta)\Delta_x \bar n_\varepsilon
-\varepsilon^2\partial_{\theta\theta} \bar n_\varepsilon
=
\bar n_\varepsilon\bigl(K(x)-\bar\rho_\varepsilon(x)\bigr),
\qquad
\bar\rho_\varepsilon(x)=\int_{\mathbb T} \bar n_\varepsilon(x,\theta)\,d\theta.
\end{equation}
As shown in \cite{PerthameSouganidis}, under standard assumptions on \(K\), \(D\), and the domain, the steady-state family admits several properties that are fundamental for the rare-mutation limit.

\begin{proposition}[Continuous asymptotic structure of the steady state]\label{prop:continuous_structure}
Assume that \(K\) is positive and nonconstant, that \(D\) is positive, periodic, and has a unique minimizer \(\theta_m\), and that the steady problem \eqref{eq:ss_n} admits a positive solution. Then the following properties hold at the continuous level.
\begin{enumerate}
\item There exists a constant \(C>0\), independent of \(\varepsilon\), such that
\begin{equation}\label{eq:rho_continuous_bound}
0\le \bar\rho_\varepsilon(x)\le C
\qquad \text{for all } x\in\Omega.
\end{equation}
Moreover, up to subsequences, \(\bar\rho_\varepsilon\) converges uniformly to a limit density \(\rho\).

\item If we write
\begin{equation}\label{eq:u_log_transform}
\bar n_\varepsilon(x,\theta)=\exp\!\left(\frac{\bar u_\varepsilon(x,\theta)}{\varepsilon}\right),
\end{equation}
then \(\bar u_\varepsilon\) is uniformly Lipschitz continuous in \(\theta\), and along subsequences
\begin{equation}\label{eq:u_limit_uniform}
\bar u_\varepsilon \to u
\qquad \text{uniformly in } (x,\theta),
\end{equation}
where the limit \(u\) is independent of \(x\), periodic in \(\theta\), and satisfies
\begin{equation}\label{eq:HJ_limit}
\begin{cases}
|\partial_\theta u(\theta)|^2 = \mathcal H\bigl(\theta,\rho(\cdot)\bigr),\\[1mm]
\displaystyle \max_{\theta\in\mathbb T} u(\theta)=0.
\end{cases}
\end{equation}

\item The effective Hamiltonian \(\mathcal H(\theta,\rho)\) is defined through the principal eigenvalue problem
\begin{equation}\label{eq:eigen_N}
\begin{cases}
-D(\theta)\Delta_x \mathcal N(x,\theta)
= \mathcal N(x,\theta)\bigl(K(x)-\rho(x)\bigr)
+ \mathcal N(x,\theta)\,\mathcal H\bigl(\theta,\rho(\cdot)\bigr), & x\in\Omega,\\
\partial_\nu \mathcal N = 0, & x\in\partial\Omega,
\end{cases}
\end{equation}
with a positive eigenfunction \(\mathcal N\). Moreover, \(\mathcal H\) has the same monotonicity in \(\theta\) as \(D\), and hence
\begin{equation}\label{eq:H_min_zero}
\min_{\theta\in\mathbb T} \mathcal H\bigl(\theta,\rho(\cdot)\bigr)
=
\mathcal H\bigl(\theta_m,\rho(\cdot)\bigr)
=0.
\end{equation}

\item The steady state concentrates on the fittest trait in the limit \(\varepsilon\to0\). More precisely,
\begin{equation}\label{eq:continuous_dirac_limit}
\bar n_\varepsilon \rightharpoonup N_m(x)\,\delta(\theta-\theta_m),
\qquad
\bar\rho_\varepsilon \to N_m(x),
\end{equation}
where \(N_m\) solves the reduced Fisher-type problem
\begin{equation}\label{eq:Nm_problem}
\begin{cases}
-D(\theta_m)\Delta_x N_m = N_m\bigl(K(x)-N_m\bigr), & x\in\Omega,\\
\partial_\nu N_m = 0, & x\in\partial\Omega.
\end{cases}
\end{equation}
\end{enumerate}
\end{proposition}

%\begin{remark}[Interpretation for the numerical method]\label{rem:AP_target_cont}
Proposition~\ref{prop:continuous_structure} provides the continuous target of the numerical scheme. The boundedness of \(\bar\rho_\varepsilon\) keeps the effective Hamiltonian in a bounded regime. The constrained Hamilton--Jacobi equation identifies the dominant trait through the maximizer of \(u\). The concentration result \eqref{eq:continuous_dirac_limit} indicates that an asymptotic-preserving discretization should recover the reduced pair \((N_m,\theta_m)\) without resolving the original density on an extremely fine trait grid. This is the sense in which the WKB reformulation is used in the sequel.
%\end{remark}


%===================================================================================================
%   Detailed scheme
%===================================================================================================

\section{WKB reformulated numerical scheme}

In this section, we develop and analyze a semi-discrete numerical scheme for
the WKB reformulation \eqref{eq:WKB_reformulation}. The scheme is formulated
in terms of the phase variable \(u_\varepsilon\) and the amplitude variable
\(W_\varepsilon\), with the total density \(n_\varepsilon\) reconstructed from
the WKB representation.

\subsection{Semi-discrete numerical discretization}

We work at the semi-discrete level in order to separate the discretization in
\((x,\theta)\) from the additional issue of time stepping. For clarity, we
present the scheme in one spatial dimension with uniform grids in the
\(x\)- and \(\theta\)-directions. This choice is not essential for the WKB
decomposition, but it keeps the notation transparent.

Let \(\{x_j\}_{j=1}^J\) be a uniform grid on \(\Omega=(a,b)\) and let
\(\{\theta_k\}_{k=1}^K\) be a uniform periodic grid on \(\mathbb T=(0,1)\),
with mesh size \(\Delta\theta\). We denote
\[
W_{j,k}^\varepsilon(t)\approx W_\varepsilon(t,x_j,\theta_k),
\qquad
u_k^\varepsilon(t)\approx u_\varepsilon(t,\theta_k).
\]
For notational convenience, we write
\[
D_k=D(\theta_k),
\qquad
K_j=K(x_j).
\]
By \eqref{eq:basic_assumptions}, we have 
\[
    0<D_{\min}\le D_k\le D_{\max}<\infty,
    \qquad
    0\le K_j\le K_{\max}.
\]


For a grid function \(V=\{V_{j,k}\}\), we use the standard second-order
difference operators
\[
\delta_x^2 V_{j,k}
=
\frac{V_{j+1,k}-2V_{j,k}+V_{j-1,k}}{(\Delta x)^2},
\qquad
\delta_\theta^2 V_{j,k}
=
\frac{V_{j,k+1}-2V_{j,k}+V_{j,k-1}}{(\Delta\theta)^2}.
\]
The homogeneous Neumann boundary condition in \(x\) is imposed by the ghost
values
\[
V_{0,k}=V_{1,k},
\qquad
V_{J+1,k}=V_{J,k},
\]
while all indices in the \(\theta\)-direction are understood periodically.
For the phase variable, we also use
\[
D_\theta^-u_k
=
\frac{u_k-u_{k-1}}{\Delta\theta},
\qquad
D_\theta^+u_k
=
\frac{u_{k+1}-u_k}{\Delta\theta}.
\]

The semi-discrete WKB system involves two numerical ingredients that are not
contained in the standard second-order difference operators introduced above.
The first one is the numerical Hamiltonian for the phase equation. The
Hamilton--Jacobi term \(-|\partial_\theta u_\varepsilon|^2\) is approximated by
\[
\widehat H(D_\theta^-u_k^\varepsilon,D_\theta^+u_k^\varepsilon).
\]
We assume that \(\widehat H\) is locally Lipschitz and consistent with the
continuous Hamiltonian, namely
\[
\widehat H(p,p)=-p^2 .
\]
We also assume that it is monotone in the sense that it is nondecreasing with
respect to its first argument and nonincreasing with respect to its second
argument. When \(\widehat H\) is differentiable, this condition reads
\[
\partial_a\widehat H(a,b)\ge0,
\qquad
\partial_b\widehat H(a,b)\le0 .
\]
Typical examples are the Godunov and Lax--Friedrichs numerical Hamiltonians,
as recalled in Appendix~\ref{append:HJ}.

The second ingredient is the conservative discretization of the transport
coupling in the amplitude equation. Recall that the continuous flux is
\[
F_\varepsilon=-W_\varepsilon\partial_\theta u_\varepsilon .
\]
We denote by
\[
\mathcal D_\theta\widehat{\mathcal F}(W^\varepsilon,u^\varepsilon)_{j,k}
\]
a conservative discretization of \(\partial_\theta F_\varepsilon\) at
\((x_j,\theta_k)\). Here \(\widehat{\mathcal F}\) approximates the flux
\(-W_\varepsilon\partial_\theta u_\varepsilon\). We assume that
\(\mathcal D_\theta\widehat{\mathcal F}\) is consistent with
\[
\partial_\theta F_\varepsilon
=
\partial_\theta(-W_\varepsilon\partial_\theta u_\varepsilon)
\]
when evaluated on smooth functions, and that it is conservative in the
periodic trait variable. For the positivity argument below, we further assume
the quasi-positivity property
\begin{equation}\label{ass:H_flux}
W_{j,k}=0,
\qquad
W_{j,\ell}\ge0\ \text{for all }\ell
\quad\Longrightarrow\quad
-\mathcal D_\theta\widehat{\mathcal F}(W,u)_{j,k}\ge0 .
\end{equation}
This property is satisfied by standard monotone finite volume fluxes, such as
the upwind and Lax--Friedrichs fluxes listed in Appendix~\ref{append:flux}.

With these two ingredients specified, we now define the discrete density and
the semi-discrete WKB system. For notational convenience, set
\begin{equation}\label{def:E}
E_k^\varepsilon(t)
=
\exp\!\left(\frac{u_k^\varepsilon(t)}{\varepsilon}\right).
\end{equation}
The reconstructed nodal density is defined by
\begin{equation}\label{def:n_reconstructed}
n_{j,k}^\varepsilon(t)
=
W_{j,k}^\varepsilon(t)E_k^\varepsilon(t).
\end{equation}


As we shall see later, the evolution of $u_\varepsilon$ and $W_\varepsilon$ 
couples with the density function. 
Here we consider a general density computed by a
trait-direction quadrature operator
\[
\mathcal Q_\theta^\varepsilon[W_j^\varepsilon,u^\varepsilon],
\]
which approximates
\[
\int_{\mathbb T}
n_\varepsilon(t,x_j,\theta)
\,d\theta 
=
\int_{\mathbb T}
W_\varepsilon(t,x_j,\theta)
\exp\!\left(\frac{u_\varepsilon(t,\theta)}{\varepsilon}\right)
\,d\theta .
\]
We define
\begin{equation}\label{eq:E_rho}
\rho_{j,Q}^\varepsilon(t)
=
\mathcal Q_\theta^\varepsilon
\bigl[W_{j,\cdot}^\varepsilon(t),u^\varepsilon(t)\bigr],
\qquad
j=1,\dots,J .
\end{equation}
A direct composite quadrature gives
\begin{equation} \label{def:m}
\rho_{j,Q}^\varepsilon(t)
=
m_j^\varepsilon(t)
:=
\Delta\theta
\sum_{k=1}^K
W_{j,k}^\varepsilon(t)E_k^\varepsilon(t),
\end{equation}
whereas in the small-\(\varepsilon\) regime one may instead use a
Laplace-type reconstruction presented in
Appendix~\ref{append:full_discretization} near the maximizer of
\(u^\varepsilon\).

Given
\[
\boldsymbol\rho_Q(t)
=
\bigl(\rho_{1,Q}^\varepsilon(t),\dots,\rho_{J,Q}^\varepsilon(t)\bigr),
\]
the discrete effective Hamiltonian
\(\mathcal H_k(\boldsymbol\rho_Q(t))\) is determined by the discrete
principal eigenvalue problem
\begin{equation}\label{eq:eigen_discrete}
-D_k\delta_x^2\mathcal N_j^{(k)}
=
\mathcal N_j^{(k)}
\bigl(K_j-\rho_{j,Q}^\varepsilon(t)\bigr)
+
\mathcal N_j^{(k)}
\mathcal H_k\bigl(\boldsymbol\rho_Q(t)\bigr),
\qquad
j=1,\dots,J .
\end{equation}
Here \(\mathcal N^{(k)}\) is the corresponding positive eigenvector, equipped
with the same discrete Neumann boundary condition in \(x\). Its normalization
does not affect \(\mathcal H_k\). For definiteness, one may impose
\[
\Delta x\sum_{j=1}^J \mathcal N_j^{(k)}=1 .
\]

With the above notation, the semi-discrete WKB scheme reads
\begin{equation}\label{eq:semi_discrete_WKB_system}
\left\{
\begin{aligned}
&\frac{\mathrm d}{\mathrm dt}u_k^\varepsilon
+
\widehat H
\left(
D_\theta^-u_k^\varepsilon,
D_\theta^+u_k^\varepsilon
\right)
-
\varepsilon\delta_\theta^2u_k^\varepsilon
=
-\mathcal H_k\bigl(\boldsymbol\rho_Q(t)\bigr),
\qquad
k=1,\dots,K,
\\[1mm]
&\varepsilon
\frac{\mathrm d}{\mathrm dt}W_{j,k}^\varepsilon
-
D_k\delta_x^2W_{j,k}^\varepsilon
-
\varepsilon^2\delta_\theta^2W_{j,k}^\varepsilon
+
2\varepsilon
\mathcal D_\theta
\widehat{\mathcal F}(W^\varepsilon,u^\varepsilon)_{j,k}
\\
&\qquad =
W_{j,k}^\varepsilon
\bigl(
K_j-\rho_{j,Q}^\varepsilon(t)
+\mathcal H_k(\boldsymbol\rho_Q(t))
-2\varepsilon\delta_\theta^2u_k^\varepsilon
\bigr),
\qquad
j=1,\dots,J,\quad k=1,\dots,K .
\end{aligned}
\right.
\end{equation}

We next show the equation satisfied by the reconstructed density
\(n_{j,k}^\varepsilon\). Since \(E_k^\varepsilon\) is independent of the
spatial index \(j\),
\[
\delta_x^2 n_{j,k}^\varepsilon
=
E_k^\varepsilon\delta_x^2W_{j,k}^\varepsilon .
\]
Moreover,
\begin{equation}\label{def:dt_n}
\varepsilon\frac{\mathrm d}{\mathrm dt}n_{j,k}^\varepsilon
=
E_k^\varepsilon
\left(
\varepsilon\frac{\mathrm d}{\mathrm dt}W_{j,k}^\varepsilon
+
W_{j,k}^\varepsilon
\frac{\mathrm d}{\mathrm dt}u_k^\varepsilon
\right).
\end{equation}
Multiplying the first equation in \eqref{eq:semi_discrete_WKB_system} by $E_k^\varepsilon W_{j,k}^\varepsilon$ and multiplying the second equation in  \eqref{eq:semi_discrete_WKB_system} by $\varepsilon E_k^\varepsilon$, 
noticing \eqref{def:dt_n},
we can derive that \(n_{j,k}^\varepsilon\) satisfies
\begin{equation}\label{eq:semi_discrete_n}
\varepsilon\frac{\mathrm d}{\mathrm dt} n_{j,k}^\varepsilon
-
D_k\delta_x^2 n_{j,k}^\varepsilon
=
\bigl(K_j-\rho_{j,Q}^\varepsilon\bigr)n_{j,k}^\varepsilon
+
R_{j,k}^\varepsilon ,
\end{equation}
where the remainder has the form 
\begin{align}
R_{j,k}^\varepsilon
&=
\varepsilon^2
E_k^\varepsilon
\delta_\theta^2 W_{j,k}^\varepsilon
-
\varepsilon
W_{j,k}^\varepsilon
E_k^\varepsilon
\delta_\theta^2 u_k^\varepsilon
-
W_{j,k}^\varepsilon
E_k^\varepsilon
\widehat H
\left(
D_\theta^-u_k^\varepsilon,
D_\theta^+u_k^\varepsilon
\right)
-
2\varepsilon
E_k^\varepsilon
\mathcal D_\theta
\widehat{\mathcal F}(W^\varepsilon,u^\varepsilon)_{j,k}.
\label{eq:R_exact}
\end{align}
Thus the reconstructed density satisfies a semi-discrete analogue of the
original density equation, with \(R_{j,k}^\varepsilon\) collecting the
trait-direction discretization effects. When the numerical Hamiltonian and
the conservative flux are evaluated on smooth exact functions, this remainder
is consistent with the trait diffusion contribution
\(\varepsilon^2\partial_{\theta\theta}n^\varepsilon\).

\paragraph{A density-compatible variant.}\mbox{}\\
The split form above is convenient because it keeps the Hamilton--Jacobi and
transport components explicit. For the stationary density estimate, however,
it is useful to record an alternative discretization in which the reconstructed
density satisfies a conservative density equation exactly. Keeping the same
phase equation, replace the amplitude equation by
\begin{equation}\label{eq:density_compatible_W}
\begin{aligned}
&\varepsilon
\frac{\mathrm d}{\mathrm dt}W_{j,k}^\varepsilon
-
D_k\delta_x^2W_{j,k}^\varepsilon
-
\varepsilon^2(E_k^\varepsilon)^{-1}
\delta_\theta^2
\bigl(W_{j,\cdot}^\varepsilon E_\cdot^\varepsilon\bigr)_k
\\
&\qquad =
W_{j,k}^\varepsilon
\bigl(
K_j-\rho_{j,Q}^\varepsilon
+\mathcal H_k(\boldsymbol\rho_Q)
+\widehat H(D_\theta^-u_k^\varepsilon,D_\theta^+u_k^\varepsilon)
-\varepsilon\delta_\theta^2u_k^\varepsilon
\bigr),
\end{aligned}
\end{equation}
where
\[
\delta_\theta^2
\bigl(W_{j,\cdot}^\varepsilon E_\cdot^\varepsilon\bigr)_k
=
\frac{
W_{j,k+1}^\varepsilon E_{k+1}^\varepsilon
-2W_{j,k}^\varepsilon E_k^\varepsilon
+W_{j,k-1}^\varepsilon E_{k-1}^\varepsilon
}{(\Delta\theta)^2} .
\]
Multiplying \eqref{eq:density_compatible_W} by \(E_k^\varepsilon\) and using
the phase equation gives
\begin{equation}\label{eq:density_compatible_n}
\varepsilon\frac{\mathrm d}{\mathrm dt} n_{j,k}^\varepsilon
-
D_k\delta_x^2 n_{j,k}^\varepsilon
-
\varepsilon^2\delta_\theta^2 n_{j,k}^\varepsilon
=
\bigl(K_j-\rho_{j,Q}^\varepsilon\bigr)n_{j,k}^\varepsilon .
\end{equation}
Thus, for this density-compatible variant, the remainder in
\eqref{eq:semi_discrete_n} is exactly
\(R_{j,k}^\varepsilon=\varepsilon^2\delta_\theta^2n_{j,k}^\varepsilon\).
Consequently, the weighted remainder can be controlled by periodic discrete
summation by parts. This observation gives a concrete way to avoid the
conditional remainder assumption used below; see
Appendix~\ref{append:remainder_condition}.



%===================================================================================================
%   Subsection: basic structural property
%===================================================================================================

\subsection{Basic structural property}

The semi-discrete scheme \eqref{eq:semi_discrete_n} can be shown to be positive-preserving 
as long as the initial data is nonnegative. 
The detailed result is summarized as the following lemma. 

\begin{lemma}[Positivity of the amplitude]
\label{lem:positivity_W}
Assume that \(W_{j,k}^\varepsilon(0)\ge0\) for all \(j,k\). Then the
semi-discrete amplitude equation in \eqref{eq:semi_discrete_WKB_system}
preserves nonnegativity, namely
\[
    W_{j,k}^\varepsilon(t)\ge0
    \qquad
    \forall\,j,k,\ t\ge0,
\]
as long as the solution exists. Consequently,
\[
    n_{j,k}^\varepsilon(t)\ge0
    \qquad
    \forall\,j,k,\ t\ge0 .
\]
\end{lemma}

\begin{proof}
It is enough to verify the quasi-positivity of the right-hand side of the
finite-dimensional ODE system for \(W^\varepsilon\). From the amplitude
equation in \eqref{eq:semi_discrete_WKB_system}, we can write
\[
\varepsilon \frac{\mathrm d}{\mathrm dt}W_{j,k}^\varepsilon
=
D_k\delta_x^2W_{j,k}^\varepsilon
+
\varepsilon^2\delta_\theta^2W_{j,k}^\varepsilon
-
2\varepsilon
\mathcal D_\theta\widehat{\mathcal F}(W^\varepsilon,u^\varepsilon)_{j,k}
+
W_{j,k}^\varepsilon B_{j,k}(t),
\]
where
\[
B_{j,k}(t)
=
K_j-\rho_{j,Q}^\varepsilon(t)
+\mathcal H_k(\boldsymbol\rho_Q(t))
-2\varepsilon\delta_\theta^2u_k^\varepsilon(t).
\]
Let \(W^\varepsilon\ge0\) and suppose that
\(W_{j,k}^\varepsilon=0\) for some \((j,k)\). Then, for an interior spatial
point,
\[
\delta_x^2W_{j,k}^\varepsilon
=
\frac{W_{j+1,k}^\varepsilon+W_{j-1,k}^\varepsilon}{(\Delta x)^2}
\ge0, 
\quad 
\delta_\theta^2W_{j,k}^\varepsilon
=
\frac{W_{j,k+1}^\varepsilon+W_{j,k-1}^\varepsilon}{(\Delta\theta)^2}
\ge0 .
\]
The same conclusion holds at the boundary because of the Neumann ghost
values in $x$-direction and the periodicity in \(\theta\)-direction. 
The quasi-positivity assumption on the conservative flux operator implies that 
\[
    -
    \mathcal D_\theta\widehat{\mathcal F}(W^\varepsilon,u^\varepsilon)_{j,k}
    \ge0 .
\]
Combining all the inequalities, noticing the fact that 
\(
    W_{j,k}^\varepsilon B_{j,k}(t)=0 
\)
since we assumed $W_{j,k}^\varepsilon=0$, 
we finally have that, at this specific $(j,k)$, the following inequality holds, 
\[
\left.
\frac{\mathrm d}{\mathrm dt}W_{j,k}^\varepsilon
\right|_{W_{j,k}^\varepsilon=0}
\ge0 .
\]
Thus the vector field points inward on the boundary of the nonnegative cone.
The standard invariance criterion for finite-dimensional ODE systems implies
the claim.
\end{proof}

%\begin{remark}[On the time-dependent stability]
%For each fixed \(\varepsilon>0\), the semi-discrete system is a
%finite-dimensional ODE system under the locally Lipschitz assumptions on the
%discrete operators. Uniform-in-\(\varepsilon\) continuation estimates for the
%time-dependent problem require a separate bootstrap argument involving the
%density, phase, amplitude, and trait-direction remainder. This issue is not
%pursued in the present work.
%\end{remark}


%===================================================================================================
%   Subsection: discrete stationary AP structure
%===================================================================================================

\subsection{Stationary AP structure under a one-sided stability condition}
\label{subsec:stationary_discrete_AP}

We now turn to stationary semi-discrete states. The goal is to identify which
parts of the continuous rare-mutation structure are retained on a fixed grid.
For the generic split WKB scheme, the only nonstandard point is the
trait-direction remainder in the reconstructed density equation. We therefore
state the stationary result under an explicit one-sided stability condition.
This keeps the AP statement conditional and avoids claiming that the WKB
remainder is automatically controlled for every trait discretization. A
verifiable density-compatible variant is discussed in
Appendix~\ref{append:remainder_condition}.

By Lemma~\ref{lem:positivity_W}, stationary states obtained from nonnegative
amplitudes satisfy
\[
    W_{j,k}^\varepsilon\ge0,
    \qquad
    n_{j,k}^\varepsilon
    =W_{j,k}^\varepsilon E_k^\varepsilon\ge0 .
\]
All estimates below are for such nonnegative reconstructed densities.

\paragraph{Weighted density balance.}
Define
\[
    \eta_j^\varepsilon
    =
    \Delta\theta
    \sum_{k=1}^K
    \frac{n_{j,k}^\varepsilon}{D_k} .
\]
Multiplying \eqref{eq:semi_discrete_n} by \(1/D_k\) and summing over the
trait variable gives the stationary weighted balance equation
\begin{equation}\label{eq:weighted_summed_density_stationary}
    -\delta_x^2m_j^\varepsilon
    =
    \eta_j^\varepsilon
    \bigl(K_j-\rho_{j,Q}^\varepsilon\bigr)
    +
    \mathcal R_j^{\varepsilon,D},
\end{equation}
where
\[
    m_j^\varepsilon
    =
    \Delta\theta\sum_{k=1}^K n_{j,k}^\varepsilon,
    \qquad
    \mathcal R_j^{\varepsilon,D}
    =
    \Delta\theta
    \sum_{k=1}^K
    \frac{R_{j,k}^\varepsilon}{D_k} .
\]
The term \(\mathcal R_j^{\varepsilon,D}\) is the discrete analogue of the
continuous weighted mutation contribution
\[
    \varepsilon^2
    \int_{\mathbb T}
    \frac1{D(\theta)}
    \partial_{\theta\theta}n^\varepsilon
    \,d\theta
    =
    \varepsilon^2
    \int_{\mathbb T}
    n^\varepsilon
    \partial_{\theta\theta}\left(\frac1{D(\theta)}\right)
    \,d\theta .
\]
At the continuous level this is bounded above by
\(C\varepsilon^2\rho^\varepsilon\). For the split WKB discretization, however,
the discrete operators act on the reconstructed density
\(n_{j,k}^\varepsilon=W_{j,k}^\varepsilon E_k^\varepsilon\), and an exact
chain rule is not available. We isolate the needed property as follows.

\begin{assumption}[One-sided weighted remainder stability]
\label{ass:remainder_stability}
There exist constants \(C_R>0\) and \(r_h\ge0\), independent of
\(\varepsilon\), such that every stationary state under consideration
satisfies
\begin{equation}\label{eq:remainder_CR_bound}
    \bigl(\mathcal R_j^{\varepsilon,D}\bigr)_+
    \le
    C_R m_j^\varepsilon+r_h,
    \qquad j=1,\ldots,J .
\end{equation}
\end{assumption}

This condition only controls the positive part of the weighted defect; a
negative contribution is favorable in the maximum-principle estimate. It is a
structural stability requirement on the trait discretization, not a consequence
of the preceding consistency statement. For a direct density-level conservative
discretization, and for the density-compatible WKB variant
\eqref{eq:density_compatible_W}, the same type of estimate follows by discrete
summation by parts. For the generic split WKB scheme, it should be checked
analytically under additional structure or monitored numerically.

We also record the mild stability condition needed for the density
reconstruction in the nonlocal reaction term.

\begin{assumption}[Reconstruction stability]
\label{ass:reconstruction_stability}
There exist constants \(c_Q,C_Q>0\) and \(e_Q\ge0\), independent of
\(\varepsilon\), such that
\begin{equation}\label{eq:quadrature_coercivity_lower}
    \rho_{j,Q}^\varepsilon
    \ge
    c_Qm_j^\varepsilon-e_Q,
    \qquad j=1,\ldots,J,
\end{equation}
\begin{equation}\label{eq:quadrature_coercivity_upper}
    \rho_{j,Q}^\varepsilon
    \le
    C_Qm_j^\varepsilon+e_Q,
    \qquad j=1,\ldots,J .
\end{equation}
\end{assumption}
For the direct composite quadrature \(\rho_{j,Q}^\varepsilon=m_j^\varepsilon\),
Assumption~\ref{ass:reconstruction_stability} holds with
\(c_Q=C_Q=1\) and \(e_Q=0\).

\begin{proposition}[Conditional density bound for stationary semi-discrete states]
\label{prop:uniform_bound_stationary_mass}
Let \((u^\varepsilon,W^\varepsilon)\) be a stationary semi-discrete state with
nonnegative reconstructed density. Assume that
Assumptions~\ref{ass:remainder_stability} and
\ref{ass:reconstruction_stability} hold. Then there exist constants
\(C_m,C_\rho>0\), independent of \(\varepsilon\), such that
\[
    0\le m_j^\varepsilon\le C_m,
    \qquad
    0\le \rho_{j,Q}^\varepsilon\le C_\rho,
    \qquad j=1,\ldots,J .
\]
\end{proposition}

\begin{proof}
Since \(n_{j,k}^\varepsilon\ge0\),
\[
    \frac1{D_{\max}}m_j^\varepsilon
    \le
    \eta_j^\varepsilon
    \le
    \frac1{D_{\min}}m_j^\varepsilon .
\]
Using \eqref{eq:weighted_summed_density_stationary} and
\eqref{eq:remainder_CR_bound}, we obtain
\[
    -\delta_x^2m_j^\varepsilon
    +
    \eta_j^\varepsilon\rho_{j,Q}^\varepsilon
    \le
    K_{\max}\eta_j^\varepsilon
    +
    C_Rm_j^\varepsilon
    +
    r_h .
\]
Let \(j_*\) be a point where
\(M^\varepsilon:=\max_jm_j^\varepsilon=m_{j_*}^\varepsilon\). Since
\(\delta_x^2m_{j_*}^\varepsilon\le0\),
\[
    \eta_{j_*}^\varepsilon\rho_{j_*,Q}^\varepsilon
    \le
    K_{\max}\eta_{j_*}^\varepsilon
    +
    C_RM^\varepsilon+r_h .
\]
If \(M^\varepsilon\le e_Q/c_Q\), the bound is immediate. Otherwise,
\(\rho_{j_*,Q}^\varepsilon\ge c_QM^\varepsilon-e_Q>0\), and hence
\[
    \frac{M^\varepsilon}{D_{\max}}
    \bigl(c_QM^\varepsilon-e_Q\bigr)
    \le
    \left(\frac{K_{\max}}{D_{\min}}+C_R\right)M^\varepsilon+r_h .
\]
This quadratic inequality gives a bound for \(M^\varepsilon\) independent of
\(\varepsilon\). The upper estimate for \(\rho_{j,Q}^\varepsilon\) then follows
from \eqref{eq:quadrature_coercivity_upper}.
\end{proof}

\paragraph{Stationary phase estimate.}
The density bound controls the coefficients in the discrete eigenvalue problem
and hence the stationary phase equation.

\begin{proposition}[Stationary phase estimate on a fixed trait grid]
\label{prop:stationary_phase_estimate}
Assume that
\[
    0\le \rho_{j,Q}^\varepsilon\le \overline\rho,
    \qquad j=1,\ldots,J,
\]
where \(\overline\rho\) is independent of \(\varepsilon\). Then the discrete
effective Hamiltonian defined by \eqref{eq:eigen_discrete} satisfies
\[
    |\mathcal H_k(\boldsymbol\rho_Q)|
    \le C_{\mathcal H},
    \qquad k=1,\ldots,K,
\]
where \(C_{\mathcal H}\) is independent of \(\varepsilon\).

Assume further that the stationary phase equation
\begin{equation}\label{eq:stationary_phase}
    \widehat H_G
    \left(
        D_\theta^-u_k^\varepsilon,
        D_\theta^+u_k^\varepsilon
    \right)
    -
    \varepsilon\delta_\theta^2u_k^\varepsilon
    =
    -\mathcal H_k(\boldsymbol\rho_Q)
\end{equation}
holds with the Godunov numerical Hamiltonian
\[
    \widehat H_G(p^-,p^+)
    =
    -\bigl((p^-)_{-}\bigr)^2
    -
    \bigl((p^+)_{+}\bigr)^2 .
\]
Then
\[
    \Delta\theta
    \sum_{k=1}^K
    |D_\theta^+u_k^\varepsilon|^2
    \le
    C_{\mathcal H} .
\]
Consequently,
\[
    \max_k |D_\theta^+u_k^\varepsilon|
    \le
    \frac{\sqrt{C_{\mathcal H}}}{\sqrt{\Delta\theta}} .
\]
If the phase is normalized by
\[
    \max_k u_k^\varepsilon=0,
\]
then
\[
    \max_k |u_k^\varepsilon|
    \le
    \sqrt{C_{\mathcal H}} .
\]
All constants are independent of \(\varepsilon\), although they may depend on
the fixed trait mesh.
\end{proposition}

\begin{proof}
The discrete eigenvalue problem is associated with the symmetric operator
\[
    -D_k\delta_x^2-(K_j-\rho_{j,Q}^\varepsilon).
\]
Its principal eigenvalue has the Rayleigh representation
\[
    \mathcal H_k(\boldsymbol\rho_Q)
    =
    \min_{\varphi\ne0}
    \frac{
        D_k\Delta x\sum_j |\delta_x^+\varphi_j|^2
        -
        \Delta x\sum_j
        (K_j-\rho_{j,Q}^\varepsilon)\varphi_j^2
    }{
        \Delta x\sum_j\varphi_j^2
    } .
\]
Since \(0\le K_j\le K_{\max}\) and
\(0\le\rho_{j,Q}^\varepsilon\le\overline\rho\), we have
\[
    \mathcal H_k(\boldsymbol\rho_Q)\ge -K_{\max} .
\]
Testing the quotient with a constant vector gives
\(\mathcal H_k(\boldsymbol\rho_Q)\le\overline\rho\). Thus
\[
    |\mathcal H_k(\boldsymbol\rho_Q)|
    \le
    \max\{K_{\max},\overline\rho\}
    =:C_{\mathcal H}.
\]

Set \(q_k^\varepsilon=D_\theta^+u_k^\varepsilon\). Then
\(D_\theta^-u_k^\varepsilon=q_{k-1}^\varepsilon\), and
\[
    -\widehat H_G
    \left(
        D_\theta^-u_k^\varepsilon,
        D_\theta^+u_k^\varepsilon
    \right)
    =
    \bigl((q_{k-1}^\varepsilon)_{-}\bigr)^2
    +
    \bigl((q_k^\varepsilon)_{+}\bigr)^2 .
\]
Summing \eqref{eq:stationary_phase} over the periodic trait grid eliminates
\(\delta_\theta^2u^\varepsilon\). Therefore
\[
\begin{aligned}
    \Delta\theta\sum_{k=1}^K |q_k^\varepsilon|^2
    &=
    -\Delta\theta\sum_{k=1}^K
    \widehat H_G
    \left(
        D_\theta^-u_k^\varepsilon,
        D_\theta^+u_k^\varepsilon
    \right)  \\
    &=
    \Delta\theta\sum_{k=1}^K
    \mathcal H_k(\boldsymbol\rho_Q)
    \le C_{\mathcal H} .
\end{aligned}
\]
The pointwise slope bound follows from the fixed-grid inequality
\(|q_k|^2\le(\Delta\theta)^{-1}\Delta\theta\sum_\ell |q_\ell|^2\). If
\(\max_k u_k^\varepsilon=0\), connecting any node to a maximizer along the
periodic grid gives
\[
    |u_k^\varepsilon|
    \le
    \Delta\theta\sum_{\ell=1}^K |D_\theta^+u_\ell^\varepsilon|
    \le
    \left(
    \Delta\theta\sum_{\ell=1}^K |D_\theta^+u_\ell^\varepsilon|^2
    \right)^{1/2}.
\]
This proves the claim.
\end{proof}

\paragraph{Formal AP trait-selection limit.}
We finally pass formally to the fixed-grid rare-mutation limit. The statement
is deliberately fixed-grid and conditional: it identifies the limiting
selection mechanism once the density and phase estimates above are available.

\begin{proposition}[Formal AP trait-selection limit]
\label{prop:formal_AP_trait_selection}
Assume that the stationary semi-discrete states satisfy
Assumptions~\ref{ass:remainder_stability} and
\ref{ass:reconstruction_stability}, and that the phase is normalized by
\[
    \max_k u_k^\varepsilon=0 .
\]
Assume also that the stationary phase estimate of
Proposition~\ref{prop:stationary_phase_estimate} applies. Then, up to a
subsequence on the fixed grids,
\[
    \boldsymbol\rho_Q^\varepsilon\to \boldsymbol\rho^0,
    \qquad
    u^\varepsilon\to u^0 .
\]
Moreover, the limiting phase satisfies the discrete constrained
Hamilton--Jacobi system
\[
    \widehat H
    \left(
        D_\theta^-u_k^0,
        D_\theta^+u_k^0
    \right)
    =
    -\mathcal H_k(\boldsymbol\rho^0),
    \qquad
    \max_k u_k^0=0 .
\]
For the Godunov Hamiltonian \(\widehat H_G\), every maximizer \(k_*\) of
\(u^0\) satisfies
\[
    \mathcal H_{k_*}(\boldsymbol\rho^0)=0 .
\]
If \(\mathcal H_k(\boldsymbol\rho^0)\) has a unique zero at \(k_m\), then
\(k_*=k_m\). Hence the selected grid trait is identified by
\[
    \theta_{k_m}\in \arg\max_k u_k^0 .
\]
\end{proposition}

\begin{proof}
Proposition~\ref{prop:uniform_bound_stationary_mass} and
\eqref{eq:quadrature_coercivity_upper} give a uniform bound on
\(\boldsymbol\rho_Q^\varepsilon\). Since the grids are fixed, a subsequence
converges to some \(\boldsymbol\rho^0\). The normalized phase is also uniformly
bounded on the fixed trait grid by
Proposition~\ref{prop:stationary_phase_estimate}; hence, up to a further
subsequence, \(u^\varepsilon\to u^0\).

Passing to the limit in the stationary phase equation gives the discrete
constrained Hamilton--Jacobi system, because
\(\varepsilon\delta_\theta^2u_k^\varepsilon\to0\) on the fixed grid. For the
Godunov Hamiltonian, \(\widehat H_G\le0\), and therefore
\(
\mathcal H_k(\boldsymbol\rho^0)=-\widehat H_G(D_\theta^-u_k^0,D_\theta^+u_k^0)
\ge0
\).
At a maximizer \(k_*\) of \(u^0\),
\[
    D_\theta^-u_{k_*}^0\ge0,
    \qquad
    D_\theta^+u_{k_*}^0\le0,
\]
so \(\widehat H_G(D_\theta^-u_{k_*}^0,D_\theta^+u_{k_*}^0)=0\). Hence
\(\mathcal H_{k_*}(\boldsymbol\rho^0)=0\). If this zero is unique, the maximizer
must be \(k_m\).
\end{proof}

The preceding results should be read as a stationary AP structure rather than
a complete AP convergence theorem. The constants are uniform in
\(\varepsilon\) for a fixed mesh, while a simultaneous mesh-refinement and
rare-mutation limit would require additional estimates on the reconstruction,
the phase slopes, and the WKB remainder.



%===================================================================================================
%	Numerical experiments
%===================================================================================================

\section{Numerical experiments}

This section specifies the numerical tests used to assess the WKB formulation.
The experiments are designed around three questions: whether the dominant trait
can be captured on a coarse trait grid, whether the method remains stable as
\(\varepsilon\) decreases, and whether the one-sided remainder condition used
in the stationary analysis is reasonable for the implemented discretization.
The numerical values in the tables below are to be filled after the production
runs.

\subsection{Common computational setting}

Unless otherwise stated, we use
\[
    \Omega=(0,1),
    \qquad
    \mathbb T=[0,1),
\]
with homogeneous Neumann boundary conditions in \(x\) and periodic boundary
conditions in \(\theta\). The baseline spatial carrying capacity and dispersal
landscape are
\begin{equation}\label{eq:numerical_KD_baseline}
    K(x)=K_0+K_1\cos(2\pi x),
    \qquad
    K_0=1,
    \quad
    K_1=0.5,
\end{equation}
\begin{equation}\label{eq:numerical_D_baseline}
    D(\theta)
    =
    D_{\min}
    +D_1\bigl(1-\cos(2\pi(\theta-\theta_m))\bigr),
    \qquad
    D_{\min}=0.2,
    \quad
    D_1=0.4,
    \quad
    \theta_m=0.35 .
\end{equation}
Thus \(D\) has a unique minimum at \(\theta_m\). In the rare-mutation limit,
this value is the expected selected trait for the continuous steady state.

The default initial data for the WKB variables are
\begin{equation}\label{eq:numerical_initial_data}
    u_k^0
    =
    -\alpha_0
    \bigl(1-\cos(2\pi(\theta_k-\theta_0))\bigr),
    \qquad
    \alpha_0=0.05,
    \quad
    \theta_0=0.70,
\end{equation}
\[
    W_{j,k}^0=1+0.1\cos(2\pi x_j),
    \qquad
    n_{j,k}^0=W_{j,k}^0\exp(u_k^0/\varepsilon).
\]
The initial maximizer \(\theta_0\) is deliberately chosen away from
\(\theta_m\) so that the computations test trait selection rather than merely
preservation of a well-prepared peak.

For the direct density computation we solve the original equation for
\(n_{j,k}\). For the WKB computation we solve
\eqref{eq:semi_discrete_WKB_system}; when explicitly stated, we also run the
density-compatible variant \eqref{eq:density_compatible_W}. The discrete
effective Hamiltonian is computed as the principal eigenvalue of
\[
    A_k(\boldsymbol\rho_Q)
    =
    -D_kL_x-	\textnormal{diag}(K_j-\rho_{j,Q}),
\]
where \(L_x\) is the Neumann finite-difference Laplacian. The phase is
normalized after each time step by the gauge transformation
\[
    c^n=\max_k u_k^n,
    \qquad
    u_k^n\leftarrow u_k^n-c^n,
    \qquad
    W_{j,k}^n\leftarrow W_{j,k}^n\exp(c^n/\varepsilon),
\]
which leaves \(n_{j,k}^n=W_{j,k}^n\exp(u_k^n/\varepsilon)\) unchanged.
Density reconstructions are evaluated in a log-stabilized form to avoid
underflow:
\[
    \rho_{j,Q}^n
    =
    \exp(s^n/\varepsilon)
    \Delta\theta\sum_{k=1}^K
    W_{j,k}^n
    \exp\left(\frac{u_k^n-s^n}{\varepsilon}\right),
    \qquad
    s^n=\max_k u_k^n .
\]

The computation is stopped when the relative stationary residual satisfies
\begin{equation}\label{eq:stationary_stopping}
    \mathcal E^n
    :=
    \max\left\{
    \frac{\|u^{n+1}-u^n\|_\infty}{\tau},
    \frac{\|W^{n+1}-W^n\|_\infty}
         {\tau(1+\|W^{n+1}\|_\infty)}
    \right\}
    \le 10^{-9} .
\end{equation}
For the direct density solver, the analogous residual is computed with
\(n^{n+1}-n^n\). In all tests the reported CPU time excludes plotting but
includes the eigenvalue solves.

We use the following diagnostics. The trait marginal and selected traits are
\[
    P_k=\Delta x\sum_{j=1}^J n_{j,k},
    \qquad
    \theta_{\rm dir}=\theta_{\arg\max_k P_k},
    \qquad
    \theta_{\rm WKB}=\theta_{\arg\max_k u_k} .
\]
The periodic distance on \(\mathbb T=[0,1)\) is
\[
    d_{\mathbb T}(a,b)
    =
    \min_{q\in\mathbb Z}|a-b+q| .
\]
The concentration width of the reconstructed trait profile is measured by
\[
    \sigma_\theta
    =
    \left(
    \frac{\sum_{k=1}^K d_{\mathbb T}(\theta_k,\theta_{\rm WKB})^2P_k}
         {\sum_{k=1}^K P_k}
    \right)^{1/2} .
\]
To monitor the assumption used in
Proposition~\ref{prop:uniform_bound_stationary_mass}, we also record
\begin{equation}\label{eq:remainder_diagnostic}
    \Gamma_R^\varepsilon
    =
    \max_{1\le j\le J}
    \frac{\bigl(\mathcal R_j^{\varepsilon,D}\bigr)_+}
         {m_j^\varepsilon+10^{-14}} .
\end{equation}
A bounded value of \(\Gamma_R^\varepsilon\) as \(\varepsilon\) decreases is not
a proof of Assumption~\ref{ass:remainder_stability}, but it is an important
numerical check that the WKB remainder is not producing an uncontrolled
positive contribution.

\subsection{Test 1: direct density solver versus WKB solver}

The first test compares the direct density discretization with the WKB solver
on the baseline data \eqref{eq:numerical_KD_baseline}--\eqref{eq:numerical_D_baseline}.
Use
\[
    J=200,
    \qquad
    K\in\{32,64,128\},
    \qquad
    \varepsilon\in\{5\times10^{-2},2\times10^{-2},10^{-2}\} .
\]
A high-resolution direct computation with \(K_{\rm ref}=1024\) is used as a
finite-\(\varepsilon\) reference when feasible; otherwise the limiting trait
\(\theta_m\) is used as the reference selection target. The quantities to
report are
\[
    e_{\rm dir}=d_{\mathbb T}(\theta_{\rm dir},\theta_{\rm ref}),
    \qquad
    e_{\rm WKB}=d_{\mathbb T}(\theta_{\rm WKB},\theta_{\rm ref}).
\]

\begin{table}[h]
\centering
\caption{Test 1. Trait-selection accuracy of the direct and WKB solvers. Fill
in the numerical values after the production runs.}
\label{tab:test1_trait_accuracy}
\begin{tabular}{cccccc}
\toprule
\(\varepsilon\) & \(K\) & \(\theta_{\rm ref}\) & \(\theta_{\rm dir}\) & \(\theta_{\rm WKB}\) & CPU ratio \\
\midrule
\(5\times10^{-2}\) & 32  & -- & -- & -- & -- \\
\(5\times10^{-2}\) & 64  & -- & -- & -- & -- \\
\(2\times10^{-2}\) & 32  & -- & -- & -- & -- \\
\(2\times10^{-2}\) & 64  & -- & -- & -- & -- \\
\(10^{-2}\)        & 64  & -- & -- & -- & -- \\
\(10^{-2}\)        & 128 & -- & -- & -- & -- \\
\bottomrule
\end{tabular}
\end{table}

The expected outcome is that, for small \(\varepsilon\), the direct method
requires a much finer trait grid to place the peak correctly, whereas the WKB
method identifies the selected trait through the smoother phase variable.

\subsection{Test 2: fixed coarse trait grid and decreasing mutation scale}

This test keeps a coarse trait mesh and decreases \(\varepsilon\):
\[
    J=200,
    \qquad
    K=64,
    \qquad
    \varepsilon\in
    \{10^{-1},5\times10^{-2},2\times10^{-2},10^{-2},5\times10^{-3}\} .
\]
The direct density solver is included until it becomes visibly under-resolved.
For the WKB solver, report the selected trait, the concentration width
\(\sigma_\theta\), and the minimum value of the discrete effective Hamiltonian
at stationarity:
\[
    H_{\min}=\min_k\mathcal H_k(\boldsymbol\rho_Q).
\]
In the rare-mutation regime, \(\theta_{\rm WKB}\) should approach
\(\theta_m\), while \(\sigma_\theta\) should shrink with \(\varepsilon\).

\begin{table}[h]
\centering
\caption{Test 2. WKB behavior on a fixed coarse trait grid as
\(\varepsilon\) decreases.}
\label{tab:test2_small_eps}
\begin{tabular}{ccccc}
\toprule
\(\varepsilon\) & \(\theta_{\rm WKB}\) & \(d_{\mathbb T}(\theta_{\rm WKB},\theta_m)\) & \(\sigma_\theta\) & \(H_{\min}\) \\
\midrule
\(10^{-1}\)        & -- & -- & -- & -- \\
\(5\times10^{-2}\) & -- & -- & -- & -- \\
\(2\times10^{-2}\) & -- & -- & -- & -- \\
\(10^{-2}\)        & -- & -- & -- & -- \\
\(5\times10^{-3}\) & -- & -- & -- & -- \\
\bottomrule
\end{tabular}
\end{table}

\subsection{Test 3: mesh dependence in the trait direction}

Here \(\varepsilon=10^{-2}\) is fixed and the trait grid is varied:
\[
    J=200,
    \qquad
    K\in\{16,32,64,128,256\} .
\]
The aim is to show that the WKB prediction of the selected trait stabilizes on
coarser grids than the direct density approximation. Plot
\(P_k\) and \(u_k\) for representative values of \(K\), and report the errors
with respect to either the high-resolution reference or \(\theta_m\).

\begin{table}[h]
\centering
\caption{Test 3. Trait-grid dependence for \(\varepsilon=10^{-2}\).}
\label{tab:test3_mesh_dependence}
\begin{tabular}{ccccc}
\toprule
\(K\) & \(e_{\rm dir}\) & \(e_{\rm WKB}\) & direct CPU & WKB CPU \\
\midrule
16  & -- & -- & -- & -- \\
32  & -- & -- & -- & -- \\
64  & -- & -- & -- & -- \\
128 & -- & -- & -- & -- \\
256 & -- & -- & -- & -- \\
\bottomrule
\end{tabular}
\end{table}

\subsection{Test 4: nontrivial heterogeneous landscapes}

To check that the method is not tuned to a single cosine profile, repeat the
stationary computation with
\begin{equation}\label{eq:numerical_KD_heterogeneous}
    K(x)=1+0.35\cos(2\pi x)+0.20\cos(4\pi x),
\end{equation}
\begin{equation}\label{eq:numerical_D_heterogeneous}
    D(\theta)
    =
    0.15
    +0.35\bigl(1-\cos(2\pi(\theta-0.25))\bigr)
    +0.08\bigl(1-\cos(4\pi(\theta-0.25))\bigr) .
\end{equation}
This dispersal landscape still has its unique minimum at \(\theta=0.25\), but
it has a sharper local shape near the selected trait. Use
\[
    J=200,
    \qquad
    K\in\{64,128\},
    \qquad
    \varepsilon\in\{2\times10^{-2},10^{-2}\} .
\]
The outputs are the total spatial density
\(\rho_j\), the trait marginal \(P_k\), the phase \(u_k\), and the selected
trait. This test should be accompanied by plots comparing the direct and WKB
reconstructions.

\subsection{Test 5: remainder and density-compatibility diagnostics}

The last test directly addresses the assumption used in the stationary
analysis. For the baseline data, compute \(\Gamma_R^\varepsilon\) from
\eqref{eq:remainder_diagnostic} for the split WKB scheme and, separately, for
the density-compatible variant \eqref{eq:density_compatible_W}. Use
\[
    J=200,
    \qquad
    K=64,
    \qquad
    \varepsilon\in
    \{5\times10^{-2},2\times10^{-2},10^{-2},5\times10^{-3}\} .
\]
For the density-compatible variant, the theory in
Appendix~\ref{append:remainder_condition} predicts a bound controlled by the
discrete second difference of \(1/D_k\). For the split scheme, the table should
be used as evidence of whether the one-sided defect remains moderate in the
computed stationary states.

\begin{table}[h]
\centering
\caption{Test 5. Numerical monitoring of the one-sided weighted remainder.}
\label{tab:test5_remainder}
\begin{tabular}{cccc}
\toprule
\(\varepsilon\) & \(\Gamma_R^\varepsilon\), split WKB & \(\Gamma_R^\varepsilon\), density-compatible WKB & selected trait \\
\midrule
\(5\times10^{-2}\) & -- & -- & -- \\
\(2\times10^{-2}\) & -- & -- & -- \\
\(10^{-2}\)        & -- & -- & -- \\
\(5\times10^{-3}\) & -- & -- & -- \\
\bottomrule
\end{tabular}
\end{table}

This diagnostic is important for presentation. If the split WKB remainder is
observed to grow rapidly as \(\varepsilon\) decreases, then the numerical
section should emphasize the density-compatible variant for stationary tests.
If it remains bounded, the table provides practical support for the conditional
stability assumption.



\bigskip
\noindent{\large\bf Acknowledgements}
X. Ruan acknowledges support from the National Natural Science Foundation of China (Grant number 12201436) and the R\&D Program of Beijing Municipal Education Commission from China, grant number KM202310028016.
W. Huang acknowledges support from .

\appendix

\section{Typical choices of the monotone numerical Hamiltonian}
\label{append:HJ}

We briefly list two choices of the numerical Hamiltonian used for the phase
equation. For the continuous Hamiltonian
\[
    H(p)=-p^2,
\]
the numerical Hamiltonian \(\widehat H(a,b)\) is required to satisfy
\(\widehat H(p,p)=H(p)\), to be locally Lipschitz, and to be nondecreasing in
its first argument and nonincreasing in its second argument.

A Lax--Friedrichs type choice is
\[
    \widehat H^{\rm LF}(a,b)
    =
    -\left(\frac{a+b}{2}\right)^2
    +
    \frac{\alpha}{2}(a-b),
\]
where \(\alpha\) is chosen larger than the maximal slope magnitude reached in
the computation.

The Godunov-type Hamiltonian used in the stationary phase estimate is
\[
    \widehat H_G(a,b)
    =
    -\bigl(a_-\bigr)^2
    -\bigl(b_+\bigr)^2,
    \qquad
    a_-:=\min\{a,0\},
    \quad
    b_+:=\max\{b,0\}.
\]
It is consistent since \(\widehat H_G(p,p)=-p^2\), and it has the required
one-sided monotonicity.

\section{Typical realizations of the discrete conservative flux}
\label{append:flux}

We describe typical realizations of the abstract conservative operator
\[
    \mathcal D_\theta\widehat{\mathcal F}(W,u)_{j,k}
\]
used in the split amplitude equation. Recall that the continuous flux is
\[
    F=-W\partial_\theta u .
\]
On a uniform finite-volume grid in the trait variable, write
\[
    \mathcal D_\theta\widehat{\mathcal F}(W,u)_{j,k}
    =
    \frac{
    \widehat F_{j,k+\frac12}
    -
    \widehat F_{j,k-\frac12}
    }{\Delta\theta} .
\]
Let
\[
    q_{k+\frac12}
    =
    \frac{u_{k+1}-u_k}{\Delta\theta},
    \qquad
    a_{k+\frac12}=-q_{k+\frac12} .
\]

\paragraph{Upwind flux.}
The standard upwind flux is
\[
    \widehat F_{j,k+\frac12}^{\rm up}
    =
    a_{k+\frac12}^+W_{j,k}
    +
    a_{k+\frac12}^-W_{j,k+1},
\]
where \(a^+=\max\{a,0\}\) and \(a^- =\min\{a,0\}\). This flux is consistent
with \(-W\partial_\theta u\) and satisfies the quasi-positivity property used
in Lemma~\ref{lem:positivity_W}.

\paragraph{Lax--Friedrichs flux.}
A Lax--Friedrichs flux is
\[
    \widehat F_{j,k+\frac12}^{\rm LF}
    =
    \frac12a_{k+\frac12}(W_{j,k}+W_{j,k+1})
    -
    \frac12\alpha_{k+\frac12}(W_{j,k+1}-W_{j,k}),
\]
with \(\alpha_{k+\frac12}\ge |a_{k+\frac12}|\). Both the upwind and
Lax--Friedrichs choices are conservative in the periodic trait variable:
\[
    \Delta\theta\sum_{k=1}^K
    \mathcal D_\theta\widehat{\mathcal F}(W,u)_{j,k}=0 .
\]

\section{A fully discrete realization}
\label{append:full_discretization}

We describe the fully discrete version used in the numerical tests. Assume
that \((u^n,W^n)\) is known at time \(t_n\). First reconstruct the total
density
\[
    \rho_j^n
    =
    \mathcal Q_\theta^\varepsilon[W_{j,\cdot}^n,u^n] .
\]
For the direct composite quadrature, use the log-stabilized formula
\[
    \rho_j^n
    =
    \exp(s^n/\varepsilon)
    \Delta\theta\sum_{k=1}^K
    W_{j,k}^n
    \exp\left(\frac{u_k^n-s^n}{\varepsilon}\right),
    \qquad
    s^n=\max_k u_k^n .
\]
A Laplace-type reconstruction may also be used as a diagnostic when
\(u^n\) has a nondegenerate maximum \(\theta_*^n\):
\[
    \rho_j^n
    \approx
    W^n(x_j,\theta_*^n)
    \exp\left(\frac{u^n(\theta_*^n)}{\varepsilon}\right)
    \sqrt{
    \frac{2\pi\varepsilon}
    {|\partial_{\theta\theta}u^n(\theta_*^n)|}
    } .
\]

For each trait node \(\theta_k\), compute \(H_k^n\) as the principal
eigenvalue of
\[
    -D_k\delta_x^2\mathcal N_j^{(k),n}
    =
    \mathcal N_j^{(k),n}(K_j-\rho_j^n)
    +
    \mathcal N_j^{(k),n}H_k^n,
    \qquad j=1,\ldots,J,
\]
with discrete Neumann boundary conditions. The eigenvector may be normalized by
\[
    \Delta x\sum_{j=1}^J\mathcal N_j^{(k),n}=1 .
\]

The phase is updated by the semi-implicit scheme
\[
    \frac{u_k^{n+1}-u_k^n}{\tau}
    +
    \widehat H
    \left(D_\theta^-u_k^n,D_\theta^+u_k^n\right)
    -
    \varepsilon\delta_\theta^2u_k^{n+1}
    =
    -H_k^n .
\]

For the split WKB scheme, the amplitude is updated by
\begin{equation}\label{eq:fully_discrete_split_W}
\begin{aligned}
\varepsilon
\frac{W_{j,k}^{n+1}-W_{j,k}^n}{\tau}
&-
D_k\delta_x^2W_{j,k}^{n+1}
-
\varepsilon^2\delta_\theta^2W_{j,k}^{n+1}
+
2\varepsilon
\mathcal D_\theta
\widehat{\mathcal F}(W^{n+1},u^{n+1})_{j,k}
\\
&=
W_{j,k}^{n+1}
\bigl(
K_j-\rho_j^n
+H_k^n
-2\varepsilon\delta_\theta^2u_k^{n+1}
\bigr).
\end{aligned}
\end{equation}
The right-hand side uses \(\rho_j^n\) and \(H_k^n\), so the update is linear in
\(W^{n+1}\) once \(u^{n+1}\) is known. Other implicit--explicit variants are
possible; the stationary tests only require convergence to the same
semi-discrete equilibrium.

For the density-compatible variant, replace \eqref{eq:fully_discrete_split_W}
by
\begin{equation}\label{eq:fully_discrete_density_compatible}
\begin{aligned}
\varepsilon
\frac{W_{j,k}^{n+1}-W_{j,k}^n}{\tau}
&-
D_k\delta_x^2W_{j,k}^{n+1}
-
\varepsilon^2(E_k^{n+1})^{-1}
\delta_\theta^2
\bigl(W_{j,\cdot}^{n+1}E_\cdot^{n+1}\bigr)_k
\\
&=
W_{j,k}^{n+1}
\bigl(
K_j-\rho_j^n
+H_k^n
+\widehat H(D_\theta^-u_k^{n+1},D_\theta^+u_k^{n+1})
-\varepsilon\delta_\theta^2u_k^{n+1}
\bigr),
\end{aligned}
\end{equation}
where \(E_k^{n+1}=\exp(u_k^{n+1}/\varepsilon)\). In implementation the ratios
\(E_{k\pm1}^{n+1}/E_k^{n+1}\) should be computed as exponentials of phase
increments, not by forming very small exponentials separately.

After each phase update, apply the gauge normalization
\[
    c^{n+1}=\max_k u_k^{n+1},
    \qquad
    u_k^{n+1}\leftarrow u_k^{n+1}-c^{n+1},
    \qquad
    W_{j,k}^{n+1}\leftarrow W_{j,k}^{n+1}\exp(c^{n+1}/\varepsilon),
\]
which preserves the reconstructed density.

\section{A sufficient condition for the one-sided remainder bound}
\label{append:remainder_condition}

We record two ways in which Assumption~\ref{ass:remainder_stability} can be
justified or checked.

\paragraph{Density-compatible discretization.}\mbox{}\\
For the density-compatible variant \eqref{eq:density_compatible_W}, the
reconstructed density satisfies \eqref{eq:density_compatible_n}. Therefore, in
the notation of \eqref{eq:semi_discrete_n},
\[
    R_{j,k}^\varepsilon
    =
    \varepsilon^2\delta_\theta^2n_{j,k}^\varepsilon .
\]
Using periodic summation by parts in the trait variable,
\[
\begin{aligned}
    \mathcal R_j^{\varepsilon,D}
    &=
    \varepsilon^2\Delta\theta
    \sum_{k=1}^K
    \frac{\delta_\theta^2n_{j,k}^\varepsilon}{D_k}  \\
    &=
    \varepsilon^2\Delta\theta
    \sum_{k=1}^K
    n_{j,k}^\varepsilon
    \delta_\theta^2\left(\frac1{D_k}\right) .
\end{aligned}
\]
Hence, for \(0<\varepsilon\le1\),
\[
    \bigl(\mathcal R_j^{\varepsilon,D}\bigr)_+
    \le
    C_{D,h}m_j^\varepsilon,
    \qquad
    C_{D,h}:=
    \max_k
    \left|
    \delta_\theta^2\left(\frac1{D_k}\right)
    \right| .
\]
This verifies \eqref{eq:remainder_CR_bound} with
\(C_R=C_{D,h}\) and \(r_h=0\). If \(D\) is smooth and the second difference is
the standard centered periodic difference, \(C_{D,h}\) remains bounded along
regular mesh refinement.

\paragraph{A sufficient condition for a generic split WKB discretization.}
For the split scheme \eqref{eq:semi_discrete_WKB_system}, the same bound is not
automatic because \(R_{j,k}^\varepsilon\) contains discrete chain-rule defects.
A simple sufficient condition is the following. Suppose that there exists
\(C_E>0\), independent of \(\varepsilon\), such that
\begin{equation}\label{eq:bounded_exp_ratio_appendix}
    \exp\left(\frac{u_{k+1}^\varepsilon-u_k^\varepsilon}{\varepsilon}\right)
    +
    \exp\left(\frac{u_{k-1}^\varepsilon-u_k^\varepsilon}{\varepsilon}\right)
    \le C_E,
    \qquad k=1,\ldots,K .
\end{equation}
Assume also that the positive part of the local reconstruction defect satisfies
\begin{equation}\label{eq:local_defect_appendix}
    \bigl(R_{j,k}^\varepsilon\bigr)_+
    \le
    C_0
    \sum_{\ell\in\{k-1,k,k+1\}}
    n_{j,\ell}^\varepsilon
    +
    r_{j,k,h},
\end{equation}
with \(C_0\) independent of \(\varepsilon\) and
\[
    \Delta\theta\sum_{k=1}^K\frac{r_{j,k,h}}{D_k}\le r_h .
\]
Then positivity of \(n^\varepsilon\), periodicity in \(k\), and the lower bound
on \(D_k\) give
\[
    \bigl(\mathcal R_j^{\varepsilon,D}\bigr)_+
    \le
    C_Rm_j^\varepsilon+r_h,
\]
for a constant \(C_R\) independent of \(\varepsilon\). Condition
\eqref{eq:bounded_exp_ratio_appendix} is stronger than a uniform bound on the
ordinary discrete derivative of \(u^\varepsilon\). Indeed,
\(|D_\theta^+u_k^\varepsilon|\le C\) only gives
\(\exp(C\Delta\theta/\varepsilon)\), which is not uniform in
\(\varepsilon\) on a fixed grid. This is why the main text treats
\eqref{eq:remainder_CR_bound} as a conditional stability assumption for the
split WKB discretization.

\section{Discretization of the original density problem}
\label{append:direct_density}

For comparison, the original density equation is discretized as
\[
    \varepsilon \partial_t n_\varepsilon
    -
    D(\theta)\Delta_{xx}n_\varepsilon
    -
    \varepsilon^2\Delta_\theta n_\varepsilon
    =
    n_\varepsilon(K(x)-\rho_\varepsilon(x)),
    \qquad
    \rho_\varepsilon(x)=\int_0^1 n_\varepsilon(x,\theta)\,d\theta .
\]
A typical semi-implicit time step is
\[
\begin{aligned}
\varepsilon
\frac{n_{j,k}^{n+1}-n_{j,k}^n}{\tau}
&-
D_k\delta_x^2n_{j,k}^{n+1}
-
\varepsilon^2\delta_\theta^2n_{j,k}^{n+1}
\\
&=
    n_{j,k}^{n+1}(K_j-\rho_j^n),
    \qquad
    \rho_j^n=\Delta\theta\sum_{k=1}^K n_{j,k}^n .
\end{aligned}
\]
The same spatial and trait grids are used for the direct and WKB solvers in
all comparison tests.

\begin{thebibliography}{99}

\bibitem{AlmeidaPerthameRuan}
L. Almeida, B. Perthame, and X. Ruan,
An asymptotic preserving scheme for capturing concentrations in
age-structured models arising in adaptive dynamics,
\emph{Journal of Computational Physics}, 464 (2022), 111335.

\bibitem{BarlesPerthame}
G. Barles and B. Perthame,
Concentrations and constrained Hamilton--Jacobi equations arising in adaptive
dynamics,
\emph{Contemporary Mathematics}, 439 (2007), 57--68.

\bibitem{MirrahimiPerthame}
S. Mirrahimi and B. Perthame,
Asymptotic analysis of a selection model with space,
\emph{Journal de Math\'{e}matiques Pures et Appliqu\'{e}es}, 104 (2015),
1108--1118.

\bibitem{PerthameSouganidis}
B. Perthame and P. E. Souganidis,
Rare mutations limit of a steady state dispersal evolution model,
\emph{Mathematical Modelling of Natural Phenomena}, 11 (2016), 154--166.

\end{thebibliography}

\end{document}
